% Generated from Review/corrected_text.md (the reviewed source text).

\bakhsh{الف. قوانین کلی و شرایط بازی}

\band{۱. برگزاری}

مسابقات جهانی اتللو (\lr{WOC}) باید هر سال برای تعیین قهرمان جهان، قهرمان زنان جهان (\lr{WWOC})، قهرمان نوجوانان جهان (\lr{YWOC}) و تیم قهرمان جهان برگزار شود. مسابقات جهانی اتللو باید تحت مسئولیت فدراسیون جهانی اتللو (\lr{WOF}) و صاحبان علامت تجاری ثبت‌شده \lr{Othello} برگزار شود. \lr{WOF} نهاد مسئول این قوانین است. تمام بازیکنان مسابقات جهانی اتللو به طور ضمنی این قوانین را می‌پذیرند. کشورهای خواستار میزبانی مسابقات جهانی اتللو باید درخواست خود را به \lr{WOF} ارائه دهند و تصمیم نهایی درباره محل برگزاری مسابقات با \lr{WOF} خواهد بود.

\band{۲. شرکت‌کنندگان}

\zirband{\LRE{۲\kern0pt.\kern0pt۱} کشورها}

مگر آنکه \lr{WOF} به گونه‌ای دیگر تصمیم بگیرد، تنها کشورهای عضو \lr{WOF} می‌توانند به مسابقات جهانی بازیکن اعزام کنند. امکان اعزام بازیکن از سوی کشورهای غیرعضو، پیش از مسابقات و توسط هیئت‌رئیسه فدراسیون جهانی اتللو تعیین می‌شود. درخواست‌های مربوط به اعزام بازیکن باید دست‌کم یک ماه پیش از مسابقات جهانی به دست \lr{WOF} برسد. \lr{WOF} دو نوع عضویت را به رسمیت می‌شناسد: فدراسیون‌های ملی و کشورهای دارای فرد رابط.

\begin{itemize}
  \item فدراسیون‌های ملی مجازند یک تیم کامل پنج‌نفره اعزام کنند (حداکثر چهار بازیکن، که بیش از سه نفر از آنان نباید از یک جنسیت باشند، به همراه یک بازیکن نوجوان اضافی).
  \item کشورهای دارای فرد رابط می‌توانند یک تیم دونفره به مسابقات اعزام کنند (یک بازیکن به همراه یک بازیکن نوجوان اضافی).
\end{itemize}

در صورتی که تعداد بازیکنان حاضر در مسابقات فرد باشد، کشور میزبان می‌تواند یک بازیکن اضافی وارد مسابقات کند (بند ب.۲ را ببینید).

\zirband{\LRE{۲\kern0pt.\kern0pt۲} بازیکنان}

\begin{itemize}
  \item تمام فدراسیون‌های ملی باید تیم ملی خود را حداکثر تا ۳۱ ژوئیه ثبت کنند. هر تیم می‌تواند یک تا چهار بازیکن (که بیش از سه نفر از آنان نباید از یک جنسیت باشند) و یک بازیکن نوجوان اضافی داشته باشد. بازیکن نوجوان باید ۱۴ سال یا کمتر سن داشته باشد.
  \item تمام کشورهای دارای فرد رابط باید تیم خود را حداکثر تا ۳۱ ژوئیه ثبت کنند. هر تیم می‌تواند یک بازیکن به همراه یک بازیکن نوجوان اضافی داشته باشد. بازیکن نوجوان باید ۱۴ سال یا کمتر سن داشته باشد.
\end{itemize}

بازیکنی نوجوان شناخته می‌شود که در ۳۱ دسامبر سال برگزاری مسابقات ۱۴ سال یا کمتر سن داشته باشد (یعنی برای مسابقات جهانی ۲۰۱۹ باید پس از ۳۱ دسامبر ۲۰۰۴ به دنیا آمده باشد). هر تغییری که پس از ۳۱ ژوئیه صورت گیرد، به اجازه صریح رئیس کل مسابقات نیاز خواهد داشت. در هر صورت، تمام تغییرات باید پیش از شروع دور اول بازی‌ها صورت گیرد. هیچ تغییری پس از آغاز بازی‌ها پذیرفته نیست. قهرمانی تیمی بر اساس سه نتیجه برتر هر تیم تعیین خواهد شد. تمام بازیکنان می‌توانند عنوان قهرمان جهان را به دست آورند. تمام زنان می‌توانند برای عنوان قهرمان زنان جهان رقابت کنند. تمام بازیکنان نوجوان می‌توانند برای عنوان قهرمان نوجوانان جهان رقابت کنند.

\zirband{\LRE{۲\kern0pt.\kern0pt۳} حداقل الزامات \lr{WOF}}

فدراسیون‌ها برای حفظ جایگاه فدراسیون خود موظف‌اند نتایج مسابقات خود را برای فهرست رده‌بندی (ریتینگ) \lr{WOF} ارسال کنند. نتایج باید تا پیش از پایان مهلت ثبت‌نام مسابقات جهانی ارسال شوند (ترجیحاً نتایج تمام مسابقات، اما دست‌کم نتایج مسابقات انتخابی برای مسابقات جهانی). برای اینکه فردی بتواند نماینده کشوری در مسابقات جهانی اتللو باشد، باید شهروند آن کشور باشد یا دست‌کم ۶ ماه از ۱۲ ماه پیش از تاریخ مسابقات جهانی اتللو را در آن کشور ساکن بوده باشد. فدراسیون‌های ملی در صورت تمایل می‌توانند قوانین سخت‌گیرانه‌تری اعمال کنند.

\band{۳. مسئولان برگزاری}

برای برگزاری مسابقات جهانی اتللو، \lr{WOF} باید خدمات یک گروه مدیریت مسابقات را فراهم کند. گروه مدیریت مسابقات شامل رئیس کل مسابقات و چند مدیر مسابقات است که هر کدام حوزه مسئولیت خود را دارند. در صورت امکان، یک نماینده از فدراسیون کشور میزبان نیز باید عضو گروه مدیریت مسابقات باشد. گروه مدیریت مسابقات مسئول آماده‌سازی و برگزاری مسابقات است؛ از جمله قرعه‌کشی‌ها و نظارت بر پایبندی بازیکنان و تماشاگران به قوانینی که در اینجا آمده است. این گروه همچنین مرجع رسیدگی به تمام اختلافاتی است که در طول مسابقات پیش می‌آید. گروه مدیریت مسابقات می‌تواند دستیارانی را تعیین کند تا در کارهای مختلف برگزاری، مانند قرعه‌کشی، داوری میزها و پخش زنده بازی‌ها، کمک کنند. در صورت امکان، باید به تعداد کافی داور میز (بند د.۲ را ببینید) برای ثبت حرکات تمام بازی‌های مسابقات وجود داشته باشد.

\band{۴. شرایط بازی}

\zirband{\LRE{۴\kern0pt.\kern0pt۱} سالن مسابقات}

کشور میزبان باید برای فراهم کردن شرایط مناسب بازی برای شرکت‌کنندگان تلاش کند. سالن مسابقات باید ساکت و به اندازه کافی روشن باشد (نور غیرمستقیم ترجیح داده می‌شود، چون نور مستقیم از سطح براق مهره‌ها منعکس می‌شود و تشخیص رنگ آن‌ها را دشوار می‌کند).

\zirband{\LRE{۴\kern0pt.\kern0pt۲} مشورت و کمک}

در طول بازی، هیچ بازیکنی نباید از هیچ یادداشتی استفاده کند، از وسایل الکترونیکی غیرمجاز استفاده کند یا هرگونه اطلاعاتی را با فرد دیگری مبادله کند (به استثنای یکی از مسئولان برگزاری مسابقات یا حریف خود، و نیز درخواست از یک تماشاگر برای فراخواندن یکی از مسئولان برگزاری به منظور حل یک اختلاف). هر مورد تخلف به رئیس کل مسابقات ارجاع داده خواهد شد تا با قضاوتی معقول درباره آن تصمیم‌گیری کند و در صورت نیاز جریمه‌ای اعمال کند.

در طول دو روز اول مسابقات (به جز بازی‌های پلی‌آف)، از زمان اعلام قرعه‌های هر دور تا پایان بازی‌شان، هیچ بازیکنی حق استفاده از وسایل الکترونیکی غیرمجاز را ندارد. برای اجرای این قانون، تمام بازیکنان باید پیش از اعلام قرعه‌ها در سالن مسابقات حضور داشته باشند.

\zirband{\LRE{۴\kern0pt.\kern0pt۳} استعمال دخانیات}

استعمال دخانیات در سالن مسابقات اکیداً ممنوع است.

\zirband{\LRE{۴\kern0pt.\kern0pt۴} وسایل الکترونیکی}

بازیکنان و تماشاگران نباید در سالن مسابقات از هیچ وسیله الکترونیکی غیرمجازی (تلفن همراه، لپ‌تاپ، ساعت هوشمند و ...) استفاده کنند. این وسایل در صورت همراه بودن باید خاموش باشند.

\zirband{\LRE{۴\kern0pt.\kern0pt۵} جمع کردن صفحه بازی}

برای جلوگیری از ایجاد سروصدا در سالن مسابقات هنگامی که بازی‌ها در جریان هستند، بازیکنان نباید پس از پایان بازی، صفحه بازی خود را جمع کنند. تمام صفحه‌های بازی باید در فاصله بین دورها یا درست پیش از شروع بازی‌ها جمع شوند و در وضعیت شروع چیده شوند.

\band{۵. تماشاگران}

حضور تماشاگران در سالن مسابقات با صلاحدید رئیس کل مسابقات مجاز است. تماشاگران و سایر بازیکنان می‌توانند بازی‌های در جریان را تماشا کنند، اما نباید به هیچ شکلی در روند بازی دخالت کنند. به ویژه، تماشاگران باید از هرگونه اظهارنظر درباره بازی یا پرت کردن حواس بازیکنان یا داور میز به هر نحو خودداری کنند (به ویژه، عکاسی با فلاش در طول بازی ممنوع است و هرگونه صدای دوربین باید غیرفعال شود). با این حال، اگر یکی از تماشاگران بی‌نظمی یا تخلفی در بازی مشاهده کند، می‌تواند آن را به رئیس کل مسابقات یا یکی از دستیارانش گزارش دهد. با صلاحدید رئیس کل مسابقات، سالن مسابقات و برخی بازی‌ها می‌توانند به صورت زنده پخش شوند (ویدئو، وب‌کم، آنلاین). رئیس کل مسابقات باید اطمینان حاصل کند که هرگونه پخش زنده یا تجهیزات مربوط به آن، باعث برهم زدن آرامش بازیکنان یا اختلال در بازی‌ها نمی‌شود.
